% Ad Familiares 11.22 --- headnote
% ad-familiares-11-22, 14 July 43 BC, Rome

\textit{Cicero to D. Brutus, from Rome on 14 July 43 BC ---
Perseus dateline \textit{Scr. Romae circ. prid. Id. Quint.
a. 711 (43)}. The subject is Appius Claudius (son of Gaius),
a young noble who had joined Antony's camp on the chance of
having his exiled father recalled. With Antony defeated at
Mutina and on the run beyond the Alps, the young Appius is
now in danger as a partisan of the wrong side, and Cicero
writes to Brutus --- by then in the field against Antony ---
to ask that he be spared.}

\textit{The letter is one of Cicero's most polished short
recommendations: the appeal turns on Brutus's already-
recognised \textit{virtus} and asks him to add \textit{clementia}
to it, while the young man's defence (pietas toward an
exiled father) is set out frankly enough that Brutus, if
need be, can borrow it as a usable pretext. It is the
second-to-last of Cicero's surviving letters to Brutus.}
